\chapter{Ingrown Toenail Surgery}

\section*{Introduction \& Clinical Indications}
Ingrown toenail surgery, formally known as partial nail avulsion with chemical matrixectomy (most commonly using phenol), is a frequently performed minor surgical procedure designed to definitively treat onychocryptosis, or an ingrown toenail. This condition predominantly affects the great toe (hallux) and arises when the lateral edge of the nail plate penetrates and grows into the surrounding soft tissue of the nail fold, leading to a cascade of inflammation, pain, and often secondary bacterial infection. The procedure involves the precise removal of the offending segment of the nail plate and, critically, the permanent destruction of the corresponding portion of the nail matrix, the specialized tissue responsible for nail growth, to prevent future recurrence. It represents a highly effective, definitive intervention aimed at restoring patient comfort, preventing recurrent infections, and improving overall foot health.

\begin{itemize}
  \item Partial Nail Avulsion with Matrixectomy
  \item Wedge Resection of Nail
  \item Phenol Matrixectomy
  \item Chemical Matrixectomy
  \item Surgical Matrixectomy (for techniques involving excision)
  \item Winograd Procedure (a specific surgical excision technique)
\end{itemize}

The fundamental surgical principle of ingrown toenail surgery is twofold: first, to immediately relieve the mechanical irritation and pressure exerted by the ingrown nail plate on the surrounding soft tissues, and second, to permanently prevent the regrowth of the offending nail segment. This is achieved by carefully separating and removing the embedded lateral portion of the nail plate, which provides instant symptomatic relief and allows the inflamed paronychial tissues to begin their healing process. The rationale for this initial step is to eliminate the direct physical cause of pain and inflammation.

The second, and most critical, principle is the permanent destruction of the underlying nail matrix cells that are responsible for producing the problematic segment of the nail. Without this step, simple nail avulsion would almost invariably lead to recurrence as the nail regrows into the same pathological pathway. Chemical matrixectomy, typically using 89\% phenol, is the most common method, leveraging its caustic properties to denature proteins and induce chemical necrosis of the matrix cells, thereby ablating them. The technical goal is to create a permanently narrowed nail plate that grows straight and freely, without impinging on the lateral nail folds, thus resolving chronic pain, preventing recurrent infections, and restoring the functional integrity of the toe.

The treatment of onychocryptosis has a long history, with early interventions often limited to simple nail avulsion, which, while providing temporary relief, was plagued by high recurrence rates. The late 19th and early 20th centuries marked a turning point with the development of more definitive surgical techniques that addressed the nail matrix. In 1929, Dr. Winograd described a procedure involving the excision of the lateral nail plate along with a portion of the underlying nail matrix and adjacent nail fold, a significant advancement in preventing recurrence. This was followed in 1937 by Dr. Zadik's introduction of a technique for total nail matrix excision, often reserved for more severe or recurrent cases.

The advent of chemical matrixectomy revolutionized the management of ingrown toenails due to its reduced invasiveness and high efficacy. In 1945, Dr. Frost pioneered the use of phenol for chemical ablation of the nail matrix. This method rapidly gained popularity, offering a simpler, less traumatic alternative to surgical excision. Over the subsequent decades, phenol matrixectomy has been refined, becoming the gold standard for permanent treatment of recurrent ingrown toenails, with ongoing research focusing on optimizing application techniques and exploring alternative chemical or physical ablative agents to further improve outcomes and minimize complications.

Ingrown toenail surgery is typically indicated when conservative management strategies have failed or when the condition presents with severe, complicated features. Specific clinical indications include:
\begin{itemize}
  \item Persistent, severe pain that significantly impairs daily activities, ambulation, and quality of life, despite diligent conservative care (e.g., warm soaks, proper nail trimming, appropriate footwear).
  \item Recurrent episodes of onychocryptosis, indicating that non-surgical measures are insufficient for long-term resolution and suggesting an underlying anatomical predisposition.
  \item Presence of chronic infection (paronychia) or recurrent acute infections, often characterized by purulent drainage, erythema, swelling, and tenderness of the nail fold.
  \item Development of significant hypertrophic granulation tissue at the lateral nail fold, which is a hallmark of chronic inflammation and often exacerbates the ingrowth by creating a moist, inflamed environment.
  \item Nail plate deformity, hypertrophy, or pincer nail deformity that predisposes to ingrowth and cannot be effectively managed by non-surgical means.
  \item Cellulitis of the toe or foot secondary to the ingrown toenail, requiring definitive source control after initial systemic antibiotic treatment.
  \item Patient preference for a permanent solution to avoid ongoing discomfort, repeated conservative treatments, and the associated lifestyle limitations.
  \item Inability to wear normal, closed-toe footwear comfortably due to chronic pain and inflammation.
\end{itemize}

Ingrown toenail surgery is generally positioned as a secondary, definitive treatment rather than a primary, first-line intervention for most cases of onychocryptosis. Initial management typically involves a trial of conservative measures, which include warm water soaks, meticulous foot hygiene, proper nail trimming techniques (cutting nails straight across, avoiding rounding the corners or cutting too short), wearing comfortable, wide-toed shoes, and, if infection is present, a course of oral antibiotics. This conservative approach is often sufficient for mild to moderate, first-time ingrown toenails.

However, for cases characterized by recurrent ingrown toenails, chronic and debilitating pain, persistent infection unresponsive to antibiotics, or significant granulation tissue formation that has failed conservative care, surgery becomes the primary definitive therapeutic procedure. In these scenarios, it is not a secondary test but rather the primary intervention chosen to achieve a permanent resolution. Furthermore, in severe acute cases, particularly in patients with underlying comorbidities such as diabetes mellitus, peripheral vascular disease, or immunosuppression, where infection control and rapid resolution are paramount, surgical intervention may be considered a primary treatment to prevent serious complications like osteomyelitis or limb-threatening infection.

\section*{Relevant Surgical Anatomy}
A thorough understanding of the anatomy of the distal phalanx and nail unit is paramount for successful ingrown toenail surgery. The nail unit comprises the nail plate (the visible hard keratin structure), the nail bed (the tissue beneath the nail plate), and the nail matrix (the germinal tissue responsible for producing the nail plate). The nail matrix extends proximally beneath the eponychium (cuticle) and is crucial for nail growth. Laterally, the nail plate is bordered by the paronychium, or lateral nail folds, which are soft tissue structures that can become inflamed and hypertrophic when impinged upon by the nail. The hyponychium is the tissue beneath the free edge of the nail.

The vascular supply to the toe is rich, originating from the dorsal and plantar digital arteries, which are terminal branches of the metatarsal arteries. These vessels form an intricate capillary network within the nail bed, essential for healing. Venous drainage mirrors the arterial supply via corresponding digital veins. Sensory innervation is provided by the dorsal and plantar digital nerves, branches of the medial and lateral plantar nerves and the deep and superficial fibular nerves, which are targeted during local anesthetic blocks. Lymphatic drainage follows the venous pathways, ultimately draining to the popliteal and inguinal lymph nodes. Surgical landmarks include the precise location of the lateral nail groove and the proximal extent of the nail matrix, which must be accurately identified for complete ablation.

\section*{Preoperative Workup \& Preparation}
Thorough preoperative workup and preparation are essential to optimize outcomes and minimize potential complications. This typically involves a comprehensive assessment of the patient's overall health. A detailed medical history is obtained, focusing on comorbidities such as diabetes mellitus, peripheral vascular disease, bleeding disorders, allergies (especially to local anesthetics), and current medications. A physical examination of the affected toe and foot assesses the extent of the ingrowth, the presence and severity of infection, vascular status (e.g., capillary refill, pedal pulses), and neurological function.

\begin{itemize}
  \item  \textbf{Infection Management:} If an active infection (paronychia or cellulitis) is present, it is often treated with warm soaks and a course of oral antibiotics for several days prior to surgery to reduce inflammation and bacterial load, though some surgeons may proceed with surgery and concurrent antibiotics.
  \item  \textbf{Informed Consent:} The procedure, its potential benefits, risks, alternatives, and possible complications are thoroughly explained to the patient, and informed consent is obtained.
  \item  \textbf{Medication Review:} Patients on anticoagulant or antiplatelet medications may require adjustment of their regimen under the guidance of their prescribing physician, though for minor procedures like this, minimal changes are often sufficient.
  \item  \textbf{NPO Status:} While typically performed under local anesthesia, if any form of conscious sedation is planned, the patient will be instructed to be NPO (nil per os, nothing by mouth) for a specified period before the procedure.
  \item  \textbf{Footwear Advice:} Patients are advised to bring loose-fitting shoes or sandals to wear home, as the toe will be bandaged and potentially sensitive.
  \item  \textbf{Hygiene:} The patient should thoroughly wash their foot and toe with soap and water on the day of the procedure.
  \item  \textbf{Prophylactic Antibiotics:} Generally not required for routine ingrown toenail surgery in healthy individuals, but may be considered in immunocompromised patients or those with significant comorbidities (e.g., prosthetic heart valves, severe diabetes) at the surgeon's discretion.
\end{itemize}

\textbf{Absolute Contraindications:}
\begin{itemize}
  \item Severe peripheral arterial disease (PAD) with critical limb ischemia: Compromised blood supply significantly impairs wound healing and dramatically increases the risk of infection, tissue necrosis, and potential limb loss.
  \item Active, spreading cellulitis or osteomyelitis: The infection must be aggressively controlled with systemic antibiotics before any surgical intervention to prevent further dissemination and worsening of the infection.
  \item Documented allergy to local anesthetics: Alternative anesthesia methods or agents must be identified and utilized.
  \item Uncontrolled bleeding diathesis or severe coagulopathy: While a minor procedure, the risk of uncontrolled bleeding is unacceptable in these conditions.
\end{itemize}

\textbf{Relative Contraindications and Precautions:}
\begin{itemize}
  \item Diabetes mellitus with poor glycemic control: Patients have an increased risk of infection, delayed wound healing, and neuropathy. Careful monitoring and optimization of blood glucose levels are essential pre- and postoperatively.
  \item Immunosuppression (e.g., HIV/AIDS, organ transplant recipients, chronic corticosteroid use): These patients have a higher risk of postoperative infection and slower wound healing. Prophylactic antibiotics may be considered.
  \item Pregnancy: While local anesthesia is generally considered safe, the use of phenol is typically avoided during pregnancy due to theoretical concerns about systemic absorption and potential fetal effects. Mechanical matrixectomy may be preferred if surgery is unavoidable.
  \item Raynaud's phenomenon or other vasospastic disorders: Caution is advised due to the potential for impaired digital circulation, which could affect healing.
  \item Active fungal infection of the nail (onychomycosis): Should ideally be treated first, as it can complicate healing, obscure underlying pathology, and potentially lead to secondary bacterial infection.
  \item Very young children or uncooperative patients: May require general anesthesia or sedation, making the procedure more complex and requiring specialized facilities.
  \item Certain dermatological conditions affecting the nail unit (e.g., psoriasis, lichen planus): May alter healing or predispose to recurrence; careful evaluation is needed.
\end{itemize}

\section*{Surgical Technique \& Procedure}
Ingrown toenail surgery is typically performed as an outpatient procedure in a clinic or office setting under local anesthesia. The general steps are as follows:

\begin{enumerate}
  \item  \textbf{Anesthesia and Patient Positioning:} The patient is positioned supine. The affected toe and surrounding foot are thoroughly cleansed with an antiseptic solution (e.g., povidone-iodine or chlorhexidine). A digital nerve block is administered at the base of the toe using a local anesthetic, such as 1\% or 2\% lidocaine (without epinephrine, especially in patients with vascular compromise), to achieve complete anesthesia of the digit. Anesthesia is confirmed by pinprick testing.
  \item  \textbf{Tourniquet Application:} A small, sterile rubber tourniquet (e.g., Penrose drain or finger tourniquet) is applied at the base of the toe to create a bloodless surgical field, which significantly improves visibility and precision during the procedure.
  \item  \textbf{Nail Plate Separation and Avulsion:} Using a sterile nail elevator or a small scalpel blade (e.g., \#15), the ingrown lateral portion of the nail plate is carefully separated from the underlying nail bed and the lateral nail fold. A straight incision is then made longitudinally through the nail plate, extending from the distal edge proximally to the nail matrix, ensuring the entire offending spicule is isolated. The separated nail spicule is grasped firmly with a hemostat and avulsed (pulled out) with a steady, upward motion.
  \item  \textbf{Matrix Ablation (Phenol Matrixectomy):} After the nail spicule is removed, the exposed nail matrix is meticulously curetted to remove any residual nail fragments, debris, or granulation tissue. A cotton-tipped applicator, soaked in 89\% phenol solution, is then applied directly to the exposed germinal matrix for approximately 60-90 seconds. This application is typically repeated 2-3 times, with a brief pause between applications, to ensure thorough chemical destruction of the matrix cells.
  \item  \textbf{Neutralization and Irrigation:} Following the final phenol application, the area is thoroughly irrigated with isopropyl alcohol or saline solution to neutralize any residual phenol and flush out necrotic tissue and debris. This step is crucial to prevent chemical burns to surrounding healthy tissue.
  \item  \textbf{Tourniquet Removal and Dressing Application:} The tourniquet is carefully removed, and hemostasis is ensured. A sterile, non-adherent dressing (e.g., petroleum jelly gauze or Xeroform) is applied to the wound, followed by a soft, bulky bandage to provide gentle compression, absorb drainage, and protect the surgical site.
\end{enumerate}
The entire procedure typically takes 15-30 minutes, and sutures are generally not required.

\section*{Complications \& Risks}
While generally considered a safe and effective procedure, ingrown toenail surgery, like any surgical intervention, carries potential risks and complications.

\begin{itemize}
  \item  \textbf{General Surgical Risks:}
    \begin{itemize}
      \item Pain, swelling, and bruising at the surgical site, which are usually mild and transient.
      \item Allergic reaction to local anesthetic agents, antiseptic solutions, or dressing materials.
      \item Bleeding, typically minor and easily controlled with direct pressure or chemical cautery.
    \end{itemize}
  \item  \textbf{Procedure-Specific Risks:}
    \begin{itemize}
      \item  \textbf{Recurrence:} The most common complication, occurring in 5-20\% of cases, often due to incomplete destruction of the nail matrix, regrowth of a small nail spicule, or inadequate removal of the offending nail edge.
      \item  \textbf{Infection:} Postoperative infection (paronychia or cellulitis) of the nail fold or wound, requiring antibiotic therapy. Rarely, more severe infections such as osteomyelitis of the distal phalanx can occur, particularly in high-risk patients.
      \item  \textbf{Nail Deformity:} The new nail may grow back narrower, thicker, discolored, or with an irregular shape (e.g., spicule formation, hypertrophic nail plate, ridging), which can be cosmetically undesirable.
      \item  \textbf{Prolonged Drainage:} Serosanguinous or seropurulent drainage from the wound can persist for several weeks, especially after phenol matrixectomy, as the necrotic tissue sloughs off.
      \item  \textbf{Granuloma Formation:} Development of hypertrophic granulation tissue at the wound site, which is a reactive inflammatory response and may require additional treatment (e.g., silver nitrate cautery, surgical excision).
      \item  \textbf{Hypertrophic Scarring or Keloid Formation:} Formation of an elevated or thickened scar along the lateral nail fold, more common in individuals predisposed to keloids.
      \item  \textbf{Nerve Damage:} Although rare, injury to small sensory digital nerves in the toe can lead to localized numbness, paresthesia, or altered sensation.
      \item  \textbf{Phenol Burn:} Accidental contact of phenol with surrounding healthy skin can cause a chemical burn, though this is minimized by careful application technique and thorough irrigation.
      \item  \textbf{Systemic Absorption of Phenol:} Extremely rare with proper technique and localized application, but theoretical risk of systemic toxicity (e.g., cardiac arrhythmias, renal damage) exists.
    \end{itemize}
\end{itemize}

\section*{Postoperative Care \& Recovery}
Immediately after ingrown toenail surgery, the toe will be bandaged, and the local anesthetic will gradually wear off over several hours. Patients are typically monitored briefly in a recovery area (PACU equivalent for minor procedures) to ensure stable vital signs and absence of immediate complications before being discharged home with detailed instructions for wound care and pain management.

During the first 24-48 hours, mild to moderate pain is expected as the anesthetic wears off, usually manageable with over-the-counter pain relievers such as ibuprofen or acetaminophen. The toe should be kept elevated as much as possible to reduce swelling and throbbing. The initial bulky dressing should be kept dry and intact, and some serosanguinous drainage is normal and anticipated. Strenuous activity, prolonged standing, and tight footwear should be avoided. Over the first 1-2 weeks, the initial dressing is typically removed after 24-48 hours, and the patient begins daily warm water or saline soaks (10-15 minutes, 2-3 times a day) to keep the wound clean and aid in the sloughing of necrotic tissue. After soaking, a fresh, non-adherent dressing (e.g., petroleum jelly gauze) is applied. The wound will continue to drain a small amount of serosanguinous fluid, which is a normal part of the healing process after phenol matrixectomy. Patients can gradually resume light activities, but high-impact sports or activities that put direct pressure on the toe should be avoided. Loose-fitting, open-toed footwear is highly recommended. Complete healing of the wound typically takes 3-6 weeks, though the area may remain sensitive for a longer period. The nail plate will eventually grow back narrower, with a straight, non-impinging edge, and patients should continue to practice good foot hygiene and proper nail care to prevent future issues.

During ingrown toenail surgery, the surgeon typically documents several key findings. These often include the presence of a hypertrophic and inflamed lateral nail fold (paronychium), significant granulation tissue formation at the site of impingement, and the precise identification of the embedded nail spicule or edge that is causing the irritation. Signs of acute or chronic infection, such as purulent drainage, erythema, and localized tenderness, are also noted. The extent of the nail plate involvement and the condition of the underlying nail bed are assessed.

While routine pathological examination of the removed nail spicule is not always performed for uncomplicated cases, if the tissue is sent to pathology, the report would typically confirm the presence of nail plate fragments and often describe inflammatory changes within the soft tissue, such as chronic inflammation, granulation tissue, and possibly signs of acute infection. If a surgical matrixectomy (excision of matrix tissue) is performed, the pathology report would confirm the presence of germinal matrix cells and potentially associated inflammatory or fibrotic changes. This pathological confirmation helps to rule out other rare conditions, though it is primarily a diagnostic confirmation rather than a guide for treatment in most routine ingrown toenail cases.

A normal and successful postoperative course following ingrown toenail surgery is characterized by a progressive resolution of symptoms and uneventful wound healing. Patients can expect a significant reduction in pain and tenderness within the first few days, allowing for a gradual return to normal ambulation and daily activities. The initial swelling and redness around the toe will subside, and any drainage from the wound will gradually decrease over 2-4 weeks as the phenol-treated tissue sloughs off and the wound epithelializes. The primary goal is for the patient to achieve a pain-free toe that allows them to wear normal footwear comfortably and resume their usual lifestyle without restrictions. The nail plate will eventually regrow, but it will be narrower and have a straight, non-impinging edge, permanently resolving the issue of ingrowth. A successful course means no signs of persistent infection, excessive pain, or recurrence of the ingrown toenail.

Diligent postoperative care and scheduled follow-up are crucial for ensuring optimal healing and preventing complications.

\begin{itemize}
  \item  \textbf{Post-operative Visits:} A follow-up appointment is typically scheduled 1-2 weeks after the procedure to assess wound healing, remove any remaining debris or slough, and provide further instructions on wound care. Additional visits may be necessary if complications arise or if healing is delayed.
  \item  \textbf{Dressing Changes and Soaks:} Patients are instructed on daily dressing changes and warm water or saline soaks (10-15 minutes, 2-3 times a day) until the wound is fully healed, usually for 2-4 weeks. Non-adherent dressings are recommended.
  \item  \textbf{Activity Resumption:} Gradual return to normal activities is encouraged, with avoidance of high-impact sports, prolonged standing, or tight footwear until the toe is fully comfortable and healed.
  \item  \textbf{Warning Signs:} Patients are educated on critical warning signs of complications, such as increasing pain, spreading redness, excessive swelling, purulent discharge, fever, or a foul odor from the wound. They are instructed to contact their healthcare provider immediately if any of these symptoms occur.
  \item  \textbf{Long-term Nail Care:} Advice on proper nail trimming techniques (cutting straight across, not too short) and appropriate footwear (wide toe box) is reinforced to prevent future ingrown toenails on other digits.
  \item  \textbf{Repeat Treatment:} In the rare event of recurrence, a repeat partial nail avulsion with matrixectomy may be considered, or an alternative matrixectomy technique might be explored.
\end{itemize}
No specific follow-up imaging or laboratory tests are typically required unless complications like osteomyelitis are suspected.

\section*{Outcomes \& Prognosis}
Ingrown toenails (onychocryptosis) are an exceedingly common foot ailment, affecting a substantial proportion of the global population at some point in their lives. It is one of the most frequent dermatological and podiatric complaints encountered in clinical practice. The incidence is notably high among adolescents and young adults, with a peak often observed in individuals between 10 and 30 years of age. While both sexes are affected, males tend to present with onychocryptosis more frequently than females. The great toe (hallux) is overwhelmingly involved, accounting for over 90\% of cases, with the lateral nail fold being the most common site of impingement. Key risk factors include improper nail trimming (e.g., cutting nails too short or rounding the corners), wearing tight or ill-fitting footwear, hyperhidrosis (excessive sweating), acute or repetitive trauma to the toe, and certain nail deformities or genetic predispositions. Although rarely life-threatening, onychocryptosis can lead to significant morbidity, including chronic pain, recurrent infections, cellulitis, and in severe or neglected cases, osteomyelitis, particularly in immunocompromised individuals or those with diabetes and peripheral vascular disease. Mortality directly attributable to ingrown toenail surgery is virtually non-existent, and severe morbidity is exceedingly rare.

The outcomes of ingrown toenail surgery, particularly partial nail avulsion with chemical matrixectomy (phenolization), are generally excellent, offering high success rates and a significant improvement in patient quality of life. Success rates for phenol matrixectomy typically range from 80\% to 95\%, indicating that the vast majority of patients experience permanent resolution of their ingrown toenail without recurrence. Surgical matrixectomy techniques also boast high success rates, often within a similar range. Functional outcomes are overwhelmingly positive, with most patients returning to normal activities, including wearing regular footwear, within a few weeks of the procedure. The primary goals of pain relief, resolution of inflammation, and prevention of infection are usually achieved effectively. Prognostic factors influencing success include the completeness of matrix destruction during the procedure, patient compliance with postoperative care instructions, and the absence of underlying medical conditions that impair healing, such as uncontrolled diabetes or severe peripheral vascular disease. Recurrence, when it occurs, is the most common reason for re-intervention and is often attributed to incomplete ablation of the matrix or the regrowth of a small nail spicule. Overall, the prognosis for patients undergoing this procedure is overwhelmingly positive, leading to lasting relief from a painful and often debilitating condition.

\section*{References}
\begin{enumerate}
  \item MedlinePlus. Ingrown toenail removal. U.S. National Library of Medicine; 2024. Available from: \url{https://medlineplus.gov/ency/article/007328.htm}
  \item Sabiston Textbook of Surgery: The Biological Basis of Modern Surgical Practice. 22nd ed. Elsevier; 2022.
  \item Schwartz's Principles of Surgery. 11th ed. McGraw-Hill Education; 2019.
  \item American Academy of Orthopaedic Surgeons. Ingrown Toenail. Available from: \url{https://orthoinfo.aaos.org/en/diseases--conditions/ingrown-toenail/}
\end{enumerate}

\clearpage
